\section{Discussion}
\label{sec:discussion}
%==========================================================================

Everything above sits inside two boundaries. These are claims about identifiability, a
property of a model and a design together, not about patients: every design number is
computed from \eqref{eq:local} with the constants of Table~\ref{tab:params}, of which
three are taken from the literature and the rest are assumed, so what the analysis establishes is
what a record \emph{of the kind assumed} could yield \emph{under the model assumed}.
Whether a real record behaves that way turns on the observation model of
Remark~\ref{sec:marker}, that a marker which is not the enzyme's substrate tracks it
monotonically, and that assumption is not testable from anything published.

The second boundary is that the model questions its own validity in the regime the
excursion argument points at --- and the boundary is narrower and better located than
``the model may not hold out there''.

\emph{Every multiple in this paragraph is a multiple of the parent's own reference level,
which is not the local model's.} The parent \eqref{eq:parent} sits at an attractor of
$\num{parent_mbar}$ at its own capacity and load, where the local model's managed level
is $\bar m = 10$; the two are different quantities and a ratio taken against one does not
transfer to the other. Read against $\num{parent_mbar}$, the parent has a clearance
maximum at $\num{peak_over_mbar}$ times it and a separatrix at $\num{sep_over_mbar}$ times
it at baseline load. The \emph{stable equilibrium is lost at the first of those, not the
second}: raising the load walks attractor and separatrix together until they collide at
the fold of \S\ref{sec:model}, and they merge at $m = \num{fold_m}$, which is the
clearance maximum, at a load of $\num{fold_load}$, or $\num{fold_load_ratio}$ times
baseline. The separatrix at $\num{sep_over_mbar}$ times the attractor is how large a
perturbation must be to escape the basin at baseline load, which is a different and also
useful number.

Both estimators here are steady-state fits, and $\num{band_beyond_pct}\%$ of a band
running to eight times the parent's reference level, $\num{band_parent}$, lies beyond that
clearance maximum, where the parent says there is no steady state to fit. Fitting the
local model to the parent's own noiseless clearance over that band recovers
$\hat\theta = \num{fit_theta_hat}$ for a patient whose true capacity is the parent
model's $\theta = 300$, against an intrinsic threshold of $\num{theta_crit}$: so
misspecification alone, with no noise and no perturbation, eats
$\num{fit_margin_eaten_pct}\%$ of a true margin of $+\num{fit_margin_true}$ and leaves
$+\num{fit_margin_est}$.

\emph{That fit has to be specified, because the number moves with its definition.} The
levels are $\num{fit_levels}$, spaced geometrically from the parent's reference level to
$\num{fit_band_top}$ times it; the load at each is the parent's own noiseless clearance
there; and what is fitted is the \emph{balance} this paper's estimator actually solves,
$s = \theta\,g(m,K_m) + k_e m + b$ with $K_m$ at its population value, $\theta$ and $k_e$
non-negative and the load constant $b$ \emph{free}, by unweighted least squares in the
load. The free $b$ is the part that matters: the estimator of \S\ref{sec:collapse}
carries a load constant it does not know, and a fit of one clearance curve onto another
does not, which makes the curve-to-curve reading a different and much more favourable
exercise.

\textbf{The fit lands on the boundary of admissibility, and that, rather than the number
it lands at, is the finding.} On every band $\hat k_e$ sits exactly on its lower bound of
zero. Released, it goes \emph{negative} --- $\num{fit_ke_hat_free}$ at this band --- and
$\hat\theta$ rises to $\num{fit_theta_hat_free_ke}$. So the local form's best
approximation to the parent's clearance here describes a patient who clears the metabolite
by a route running backwards, and the best \emph{admissible} approximation describes one
with no non-enzymatic clearance at all. Neither is a patient. \emph{Over this band the
local family represents the parent only at its own edge.}

Two things follow, and the first disarms the number. \textbf{A boundary solution is not a
point estimate}: $\hat\theta = \num{fit_theta_hat}$ is where the constraint stops the fit,
not where the likelihood peaks, so the gap between it and
$\theta_{\mathrm{crit}} = \num{theta_crit}$ is not a margin and should not be read as
one --- a conclusion that holds whatever the spacing and the weighting, and therefore more
firmly than the spread below. \textbf{And nothing in the fit would reveal any of it.} The
admissible fit's residual over these levels is $\num{fit_resid_sigma}$ of a single $5\%$
error in the recorded load, which is the accuracy
Section~\ref{sec:tierload} asks for: the misspecification is invisible in exactly the
diagnostic that would be used to look for it, which is the same \emph{symptom}
Section~\ref{sec:structuredep} finds for the published alternative by a different
mechanism --- \S\ref{sec:pattern2} takes the two together.

\emph{What is robust here is that the margin is consumed, not what is left of it.} Three
equally defensible readings move $\hat\theta$ at the same band: weighting the residual by
the load, which is what a multiplicative load error implies, gives
$\num{fit_theta_hat_rel}$; reading ``over the band'' as a dense grid rather than the
design's own nine levels gives $\num{fit_theta_hat_dense}$; spacing the nine levels
linearly gives $\num{fit_theta_hat_lin}$. Against
$\theta_{\mathrm{crit}} = \num{theta_crit}$ those span an estimated margin from about
$-4$ to $+9$, so \textbf{the sign of what survives is not determined by the data at all},
and the $+\num{fit_margin_est}$ above is the design's own nine levels and nothing more
general. The conclusion that carries is the one that holds across every reading: a
noiseless, unperturbed fit of the admissible local form to the parent over this band
consumes essentially the whole of a $+\num{fit_margin_true}$ margin. And the two models
disagree by a factor of $\num{clear_ratio}$ about the largest load a patient can clear:
the local model needs $\num{hold8_ratio}$ times baseline to hold eight times \emph{its}
managed level, $\num{band_local}$, while the parent cannot clear more than
$\num{fold_load_ratio}$ times baseline at all.

\textbf{What the boundary does \emph{not} reach is the design's own operating points, and
that is a bound rather than a hope.} The local model's marker obeys
$\dot m \le L - k_e m$, so every trajectory at the recorded load is capped at
$L/k_e = \num{local_ceiling}$; and the parent's clearance maximum cannot fall below
$\sqrt{K_c M_i} = \num{peak_floor}$ at any capacity, because that is where the alternative
route's own derivative turns. A margin of $\num{peak_margin}$, so \emph{no} operating
point of any design in this paper lies past the parent's clearance peak, at any capacity
in Table~\ref{tab:scope}. Only the excursion band --- which the paper declines to use ---
goes beyond it. That margin is a property of the constants carried here and not a theorem:
across this paper's own prior range for $K_c$, $M_i$ and $\rho_0$, fewer than
$\num{prior_floor_pct}\%$ of parameter draws even place $\sqrt{K_c M_i}$ below
$L/k_e$, and fewer still actually reverse it.

The decompensation route is therefore excluded three times over: the excursion it would
supply is not observed at the magnitude required, the model that would interpret it does
not hold there, and the reduction through which that model is evaluated is not reliable
there either.

The third reason has the same shape as the first two, and it is the one this paper acquired
last. Everything computed from \eqref{eq:alphaqss} anywhere above is computed at a fixed
point, and at a fixed point the reduction is exact as an identity rather than as an
approximation (Section~\ref{sec:reduction}). A decompensation is precisely the regime in
which the quantity of interest is a trajectory and not a fixed point, and there the two
formulations part company: the reduced and two-state systems differ by $3.4\%$ to $25.4\%$
along the approach to equilibrium, and the basin boundary itself sits up to $49.6\%$ away
from the reduced separatrix when carnitine starts at its pre-perturbation value, which is
what a sudden load spike leaves behind. The displacement is signed and the unsafe sign is
the reachable one: from a depleted pool the boundary falls \emph{below} the reduced
separatrix, so there the closed form calls safe a state the parent escapes from. So the one
regime that would supply the excursion is also the one regime in which the closed form
carries no warrant. The route is not narrowly missed; it is excluded on the excursion, on
the model, and on the reduction, independently.

\subsection{What the five results amount to}

\textbf{The contribution is the design rules, with a scope condition attached, and not the
capacity estimate.} What transfers is a set of statements about identifying a saturable
channel from a marker under a known periodic input: carry the input as a regressor, never
sample at its period or a multiple of it, record the driving load as eaten and to a stated
accuracy, and establish where in the parameter range the target is recoverable before
quoting a precision anywhere in it.

The first result is actionable now and survives the scope limit: the excitation is already
being administered, and a protocol loses a factor of $\num{alias_cost}$ against monthly draws by scheduling bloods
on dosing days, which is still a factor of $16$ at a disease-range capacity where the estimate itself
is gone. The second should change how the problem is posed. Near normal capacity the
decision is blocked not by identifiability but by its own boundary, which is stipulated
rather than measured and which no better estimate of $\theta$ and no larger cohort can
repair; in the disease range both block. The third and fourth bound what should replace it, and the fifth says what the whole of it
is relative to.
The only endpoint evidence reports a smooth dose--response and no threshold, which argues
against the dichotomy rather than for a better-placed one, and an expected event rate does
not serve either: at the operating point the model-based rate is thirteen times
\emph{worse} than reading the same rate off the measured marker, the published exponent is
known only to a factor of $4.8$, and the map has no published reference concentration to
anchor it.

What is left where capacity is not recoverable is not nothing: $K_m$ is determined at every
capacity, the record still supplies an upper bound on $\theta$, and every design rule
holds. What is gone is the point estimate, and that argues for an external measurement
rather than a larger study.

\subsection{A common pattern across three of them}

Three degeneracies appear in this paper, reached by unrelated routes and concerning
different parameters. They have the same shape, and naming it is the most portable thing
here. The excitation is a product of a therapy factor and a marker factor,
$-(1 + \Delta(t)/\theta)\,\varphi(m)$, of which the literature had measured only the
second, so what the dose train contributes is not a bigger increment but a \emph{moving}
one. Proposition~\ref{prop:rank} is exactly singular at a constant recorded load and names
its own cure, so what closes it is the day-to-day \emph{variability} of ordinary eating,
while the record's \emph{accuracy} is a separate requirement entering through $\sigma$. And
the channel-exchange result (supplement, \S S1.11.1) breaks an exchange only when the inputs on
two channels differ \emph{non-constantly}, a constant increment being absorbed into the
amplitude it rides on however large it is.

In every case the separating quantity is a difference that moves, and in every case the
intuitive alternative fails: a bigger increment does not substitute for a varying one, nor
a more precise load record, nor a larger constant input. \emph{Magnitude, precision and
presence are each the wrong axis, and each is the axis a study would naturally be designed
to improve.} Two of the three are worse than neutral. A dietary record \emph{tightened} by
averaging to a weekly mean is more precise and destroys the variability
Proposition~\ref{prop:rank} needs, so Section~\ref{sec:spec} asks for a data-cleaning step
to be omitted rather than a measurement added; a \emph{steadier} administration route
removes the therapy's excitation entirely (Section~\ref{sec:route}). A larger dose is the
exception, since it raises $\Delta/\theta$ and does help, which is worth saying because it
is the only one of the three obvious moves that does.

This is a pattern across three results and not a theorem: they are three calculations that
happen to resolve the same way, and nothing here says they must. What can be said is
narrower and still useful. In a saturating compartment model observed through one marker, a
design's information sits in what varies across its samples rather than in the levels the
samples sit at, and a reader carrying that to another system should check the three
separately rather than assume the shape transfers.

\subsection{A second pattern, across the two misspecification results}
\label{sec:pattern2}

Two results here are about a \emph{wrong structure} rather than a degenerate one, and they
were reached by routes with nothing in common. Section~\ref{sec:structuredep} fits
\eqref{eq:local} to records generated by a published alternative readout and recovers a
capacity too \emph{high}, by a factor of $1.3$ to $2.0$ at this paper's own dose. The
opening of this section fits the same form to the parent's own noiseless clearance across
the excursion band and recovers one too \emph{low}, $\num{fit_theta_hat}$ against a true
$300$. The mechanisms are unrelated: the first is a free marker offset carried alongside a
free amplitude, the second a curved alternative channel that a linear $k_e m$ cannot
follow, which drives $\hat k_e$ onto its bound.

What they share is the symptom, and that is the portable part. \textbf{In both, the fit's
residual stays inside the noise the specification already asks for} --- under a fifth of
one $\sigma$ in the first, $\num{fit_resid_sigma}$ of a single load error in the second
--- \textbf{while the capacity it returns is wrong.} So on this model \emph{residual size
does not bound structural error}: the diagnostic a reader would naturally reach for to
detect a wrong structure is the one diagnostic that cannot see it. That is why
Section~\ref{sec:tierstruct}'s single blood draw is a \emph{structure} test and not a
goodness-of-fit threshold --- it separates the two structures where they disagree, in the
minutes after an infusion, instead of asking how well either fits.

And the two err in \emph{opposite} directions, which forecloses the obvious repair. A bias
of known sign could be carried as a correction or as a one-sided bound; these cannot,
because nothing in a record says which structure is the wrong one. As above this is a
pattern across two results and not a theorem. What can be said is narrower and still
useful: a small residual is evidence about noise and not about structure, and a study that
means to establish a structure has to do it with a designed discriminant rather than with a
fit.

\subsection{What would change the answer}

A same-day paired repeatability figure for the breath test, which does not exist and which
Tier~2 of Section~\ref{sec:spec} asks for, since everything the therapy channel buys is
conditional on the delivered increment being characterisable. An endpoint model that does
not require a threshold, built on the natural history study's own hazard formulation, which
needs the reference concentration and the model specification that publication does not
report. A marker that is the substrate, or an assay of it, which would remove the
observation-model assumption the whole construction rests on and is the single change that
would most enlarge what a record can support. And an external anchor on $\theta$: capacity
is identified relative to the model structure, and Section~\ref{sec:structuredep} exhibits
a published alternative under which it is not identifiable at all, fitted at a residual too
small to notice, while Section~\ref{sec:scopeprec} reaches the same place from the flux
identity and adds that in
the disease range the record does not bound $\theta$ from below whatever the structure.
Tier~2 is the only anchor here outside the longitudinal record, and a second one
independent of it would turn every statement about the value of $\theta$ from conditional
into checkable.

\vspace{0.5em}
\hrule
\vspace{1em}

%==========================================================================
